% =====================================================================
%  R. Nielsen, AN INTERPRETATION OF THE SACRED HISTORY
%  In: For Idea and Reality.  A Journal, edited by R. Nielsen, B. Bjørnson and Rud. Schmidt.  Year 1870. First Volume (Bind 2).
%  Copenhagen: Chr. Steen & Søn, Publishers; Printed by Louis Klein.  Bind 2, pp. 289-324.  (En Fortolkning af den hellige Historie, in For Ide og Virkelighed.)
%  ENGLISH TRANSLATION, from transcription.tex (Danish), the text of record.
%  GENERATED by ebook-method/for-idee-1870-73/tr_build.py --write from this template and the translators' parts in
%  ebook-method/for-idee-1870-73/tr/ (TRANSLATOR-BRIEF.md; terms: GLOSSARY.md, its key terms approved by Hans
%  2026-10-13/14).  Conventions: TRANSLATION-PLAYBOOK.md.  Page numbers of the print in the margin.
% =====================================================================
\documentclass[12pt,a4paper,openany]{book}
\usepackage[utf8]{inputenc}
\usepackage[T1]{fontenc}
\usepackage{libertinus}
\usepackage{libertinust1math}
\usepackage{textalpha}
\usepackage[english]{babel}
\usepackage[protrusion=true,expansion=true]{microtype}
\usepackage{setspace}
\usepackage[margin=1.3in]{geometry}
\usepackage{fancyhdr}
\usepackage{hyperref}
\hypersetup{pdftitle={An Interpretation of the Sacred History}, pdfauthor={Rasmus Nielsen}, hidelinks}
\setlength{\headheight}{15pt}
\pagestyle{fancy}
\fancyhf{}
\fancyhead[LE,RO]{\thepage}
\fancyhead[C]{\textit{R. Nielsen: An Interpretation of the Sacred History (1870)}}
\setstretch{1.3}
\usepackage{marginnote}
\renewcommand{\thefootnote}{\arabic{footnote})}   % the print numbers each page's notes 1), 2)
\newcommand{\opage}[1]{\marginnote{\footnotesize\textit{[#1]}}}
\newcommand{\essayhead}[2]{\begin{center}{\large #1}\ifx&#2&\else\\[0.3em]{\normalsize #2}\fi\\[0.5em]\rule{3em}{0.4pt}\end{center}\medskip\noindent}
\newcommand{\signature}[1]{\par\smallskip\hfill #1\hspace*{2em}\par}

\begin{document}

% the title page of Bind 2 (1870-1-for-idee-og-virkelighed.pdf), in English
\begin{titlepage}
\begin{center}
\vspace*{3em}
{\Huge For Idea and Reality.}\\[1em]
\rule{8em}{0.4pt}\\[1em]
{\large A Journal,}\\[0.8em]
{\small edited}\\[0.4em]
{\footnotesize by}\\[0.6em]
{\large R. Nielsen,}\\[0.4em]
{\large B. Bjørnson and Rud. Schmidt.}\\[1em]
\rule{8em}{0.4pt}\\[0.6em]
{\large Year 1870. First Volume.}\\[0.6em]
\rule{8em}{0.4pt}\\[4em]
Copenhagen.\\[0.4em]
Chr. Steen \& Søn, Publishers.\\[1.5em]
{\small Printed by Louis Klein.}\\[3em]
{\small Translated from the Danish}
\end{center}
\end{titlepage}

\essayhead{An Interpretation of the Sacred History.}{}
% --- p. 289 ---
\setcounter{footnote}{0}%
\opage{289}When great questions are in motion, and the views fit to bear arms
take the field, it is not unusual
for the campaign to be opened with a sharp declaration of war.
Each of the contending parties then regards
its standpoint as a rock-solid stronghold, an impregnable
fortress, from which sorties can be made in various directions
against the main force of the enemy, while the standpoint
itself is secured on every side against any attack.
Here there is no talk of negotiations, amicable
agreements or mutual concessions. The whole
truth is on the one side, the whole untruth on
the other; so it is \emph{to be or not to be}: what
cannot stand must fall. Worldly-wise prudence
does not like decisions of this kind; it assumes, without having looked,
that there is right and wrong on both sides. So
one hushes the combatants and would gladly conjure away
the storm. But it is in vain, when breathing is
heavy and the air close.

In the life of culture there are intermediate states, when it almost
looks as if mankind, out of sheer cleverness,
were about to become foolish, when one explains
% --- p. 290 ---
\setcounter{footnote}{0}%
\opage{290}the phenomena without troubling about their essence
and can answer all questions about them save those that actually
concern the matter. The powers of the spirit slacken; reason dozes,
one would think it were taking an afternoon nap, while
a thousand thoughts of waking understanding swarm in the day
like humming gnats in the heat. It is in such intermediate states
that the doctrine of the middle way flourishes, for ``the middle way
is broad'' and grows ever broader. A pretended
critical, but in reality uncritical, heaping together
of all manner of literary remains, an indigestible mixture of
heterogeneous components --- fragments of old, shattered,
collapsed theories --- accumulate and settle
fast as a dam against all principles, like stirred-up
mud against the fresh current of thoughts, like dregs and gravel
that clog the wellspring of the ideas. Under such circumstances
it will then be necessary that it come to
a break, to a final reckoning, and that the challenge be to
life and death.

But a dispute can also be conducted on peaceful
terms. Instead of beginning with an emphasis
on the points that exclusively mark the disagreement,
one may reverse the procedure and make a beginning
by bringing out the presuppositions on which
complete agreement prevails. With the agreement as starting point
it will then gradually be possible to make it evident
in what the disagreement consists, whence it springs,
and what significance it actually has. In the present
case, where the subject is an interpretation of the sacred
history and of the criticism that forms the foundation of the interpretation,
we choose the last-mentioned procedure.
We do so all the more gladly because the work to which
% --- p. 291 ---
\setcounter{footnote}{0}%
\opage{291}we in this treatise intend to refer the judgment of
the credibility of the evangelical accounts is a fruit
of many years of careful research, a work of scientific
value, which presents essential points of agreement,
namely: \emph{Interpretation} of the Synoptic \emph{Gospels
by Dr. H. N. Clausen} (Copenhagen 1850) (Fortolkning af de synoptiske Evangelier).

In the part of the Introduction that deals with ``the
historical character of the Gospels'' the reverend interpreter,
as far as the scientific side of criticism is concerned, has set up
principles which, in their generality, must indeed
be regarded as unassailable. We can summarize
these principles under the following three points of view. 1) \emph{The ancient ecclesiastical assumption that in}
``\emph{the subject of the evangelical account,
Christ's earthly life and ministry, one had a
purely historical material, a series of indisputable
facts}'', \emph{is critically untenable.}
A distinction must be made between profane and sacred history.
``The historical material that fills the evangelical account
cannot, with regard to its inner nature,
be likened to anything else\dots{} It is the wonderful, the
inexplicable, the supernatural that is here the account's
red thread, that indicates its keynote; even where
the miracle does not step forth immediately, where the events
for a time fit into the accustomed course of the
natural connection, the whole nevertheless moves on the dark,
holy ground of mystery''. The interpretation must therefore be determined
by the conception of the miracle.

2) \emph{To deny the possibility of miracle is to deny
the sacred history.} Yet it by no means follows from this
that one should be minded to
% --- p. 292 ---
\setcounter{footnote}{0}%
\opage{292}give up the possession of such a history, ``in so far as
by this is understood a history of the wisest teacher in the world
and the holiest man, or the possession of a Holy Scripture
in which this man's teaching and the history of his life have been
preserved''. But the task becomes that of explaining the
supernatural in a natural way; and to solve this
task one can go two ways.

One can, on the presupposition that the narrative
contains a real occurrence, a truthful event,
in which everything took place naturally, seek to explain
the semblance of the supernatural partly by the evangelist's
lack of ability for clear and distinct presentation, partly
by his limited conception of the connection of things
and his original bent to see the natural
in a supernatural light. Thus, for example, Jesus' words to
Peter on the occasion of the demanded temple tax (Mtth.
17:27): ``Go thou to the sea, and cast a hook and take the
first fish that cometh up; and when thou openest its
mouth, thou shalt find a stater; take that and give it to them
for me and thee'', are, according to Dr. Paulus, supposed to be a speaking
proof of the evangelist's linguistic helplessness; for
what he actually wants to tell is that Christ said to
Peter: ``Thou art a fisherman, cast the fishing-hook, and from the
first best fish that thou loosest from the hook thou wilt
doubtless raise the money needed''. The historical
ground on which the evangelical account must be assumed
to rest is, under this treatment, more and more hollowed out,
the kernel of full reality dwindles away,
and only the shell remains.

So one came then --- ``at an even pace'' --- onto
the second way, ``the mythicizing treatment'', whose
% --- p. 293 ---
\setcounter{footnote}{0}%
\opage{293}most significant representative is Strauss. That the sacred
history is a myth, a web of myths, means that
what we have before us here are not real facts, but only religious
representations; ``Thought works
after the manner of instinct. It is the strong, as yet undeveloped
urge to clarify for itself what fills the innermost
being of the soul that sets the powers of the spirit in motion;
it is the moral-religious ideas that work their way out of
the spiritual state of ferment toward clarity and rounded
definiteness, as they step forth in bodily form
as definite personalities, in a manifold of
figures, positions in life and deeds''. The sacred history
is, in other words, a product of unconsciously poetizing
imagination.

3) \emph{The basic condition for a true historical conception
of the sacred history is a believing
acknowledgment of the possibility of miracle.} ``That a
history which, precisely as historical reality, has become
the turning point in the religious history of the world, and in the
long series of centuries has passed over and into
the Christian life, should let itself be dissolved and spun out into
an idealistic-speculative symbolism --- that the personal
life of Christ, which in all its main stages of humiliation and
exaltation, of suffering and glorification has become the center
of a reality of the highest significance, of
a communal faith and a communal life with world-transforming
power, of devout contemplation and worship,
of poetry and art, of praise and imitation,
should all at once be torn out of the bright day of clear
reality and set in a hazy twilight, since by
an unexampled and inexplicable illusion it is supposed to have happened,
% --- p. 294 ---
\setcounter{footnote}{0}%
\opage{294}a series of symbolizing narratives --- one knows not
when or how --- came to be accepted as
real history of a teaching and miracle-working, a
suffering and crucified, a risen and ascended Redeemer
and Savior, an object of confession and invocation
in an ever greater sphere: this is so monstrous
a claim that the Christian world, not excepting even its great
indifferent part, would have to be brought to pause,
to stop thoughtfully and enter into itself''. The
sacred history is therefore the history of miracle, a real
history of real miraculous deeds.

On these points there is here, as said, complete
agreement; but now there is in the evangelical accounts
not merely something that is an object of scientific
examination; here there is also something that is an object of
faith: the history, as sacred history, is a history of faith.
When the relation between faith and knowledge is now to be more closely
determined, the disagreement begins.

That faith and knowledge must, as it is put, be reconciled
with one another, is granted by all who assume that religion
and science can be united in one consciousness.
It is not the possibility of reconciliation, but its conditions,
about which the dispute is.

Do faith and knowledge have the same principle, or are there here two
mutually different, absolutely heterogeneous principles? That
is the question. The reverend interpreter assumes
the former, we the latter; an invaluable aid
to clarifying the matter we have in the sacred history.

If the scientific use the author has
made of the principles he has set up can really furnish
proof of the identity of principle of faith with knowledge, then we see
% --- p. 295 ---
\setcounter{footnote}{0}%
\opage{295}ourselves compelled to give up our way of looking at things; but if it turns
out that the interpretation before us speaks precisely in the
most forcible way for a difference in principle between faith
and knowledge, then the interpreter himself surely cannot deny
such a difference without coming into conflict with his
own principles.

Scientific criticism is, as is also especially
stressed, to aim at ``a separation of superior
and subordinate parts of Scripture, more or less
warranted in the historical respect''. ``When the Gospels thus
contain, alongside accounts whose origin
can be traced back to the apostolic eyewitness, a few
narratives which are distinguished from the properly historical
kernel of the Gospels by the complete uncertainty concerning
the paths of transmission by which they have come
into the hands of the evangelical authors (so the evangelical
pre-history in the first chapters of Matthew
and Luke), when, over against accounts in which what happened
and is told steps forth in clear visibility,
there stand a few others that betray an uncertain, doubting conception
in the narrating witness himself, and thus let
the event told remain in a mystic half-darkness
(this holds, for example, of the temptation in the wilderness, of the transfiguration
on the mountain, of the walking on the sea, of
the resurrection of the dead of Jerusalem); when, alongside
many merely apparent discrepancies, where
the semblance of mutual contradiction dissolves on closer
consideration into a higher or lesser degree of accuracy
and completeness, there are found a few differences
that cannot be removed in any satisfactory
way (as, for example, the indication of the time of the last
% --- p. 296 ---
\setcounter{footnote}{0}%
\opage{296}Passover meal); then it cannot be the right way --- because
it is a way that leads to no resting-point ---
when a concept of the character of divinity is here set up
which by a decree of power is to strike down doubt and
frighten objection back. For it is to put
the cause of faith on the extremest and most dangerous point to
demand the unnatural of the believer: that he is to keep
his eyes tightly shut, or that, seeing, he is to renounce
the ability to see'':

This way of looking at things must, in the scientific sense,
be regarded as entirely justified. The reverend
interpreter, in pleading the right of criticism, is so far
from making too strict demands on science that
we must rather regret that in carrying out these
demands he has not been strict enough. But what follows
now from this with regard to faith?

It is remarkable, however, that precisely those sections which, according to
the above statement, belong among the historically uncertain, not
to say unhistorical, are, with few exceptions,
included among those which in the course of the church year are read aloud
for the edification of the congregation. If we now hold fast with
the author to the presupposition that the truth of faith
in the sacred history is to be clarified, grounded and
determined by scientific interpretation of Scripture, then it is
easy to see what follows from this. For from this it follows: that
\emph{in the proportion that the credibility of those evangelical accounts is
weakened in science,} must
\emph{it also be weakened in the congregation.}

If science is really to exert any influence in principle
on faith, then it must obviously be
able, little by little, to transform the congregation's
% --- p. 297 ---
\setcounter{footnote}{0}%
\opage{297}religious consciousness, make it attentive to the nature
of ``the historical material in the Gospels'', open
its eye to ``a separation of superior and subordinate parts,
more or less warranted in the historical respect''. As long as the congregation continues to listen
to the Gospel of the Annunciation, to the Gospel of Christ's
birth, of the Wise Men from the East who saw the star, etc.,
and mark well, listens without scientific criticism, it
will naturally take the events quite as they
are told, ascribe to the supernatural the same historical
reality as to the natural, and think that it is all
--- ``both certain and true in every way''. But
as the hearers, with the help of science,
are put deeper into the critical view, faith
and with it hearing will change. It will then be understood
that these ``fragmentary narratives'' are of
doubtful origin, that ``a mythicizing dressing'',
an ``embellishing setting'', has stepped in ``supplementing and
shaping'' to make good ``the scantiness and wavering''
in the historical traditions. The congregation will
then, when it hears the angels' song of praise, not let itself be
infatuated by a conception of real angels, but
content itself with the real shepherds --- for the reality
of the shepherds scientific criticism can vouch ---
but as for the angels, will bethink itself of ``the dramatic
use of the Jewish conceptions of the appearance of the world of spirits
in human fashion in the human
world, in Jewish form, even to the giving of names
and ranks''. The enlightened congregation will in faith
understand that the miracle of the star is not any proper
miracle, but a remarkable phenomenon, composed of
% --- p. 298 ---
\setcounter{footnote}{0}%
\opage{298}a natural sign in the heavens and an optical illusion; ``for if
I can say of a star that it goes before me, while
rather I go in the direction toward it, then
I can also say that it comes to stand here
or there, while it is rather I who come
to stand in such a way that I see it just there, as
over a house''.

From the scientific standpoint of faith it has often been
complained that the personal devil is to be in
the history, and that it is extremely difficult for science
to get him driven out of the consciousness of believers. An
essential obstacle lies without doubt in the fact that the congregation
every year on the first Sunday in Lent hears the Gospel
(Mtth. 4) about the temptation in the wilderness, and --- naturally
without any criticism --- takes the narrative as a real
event, an essential part of the sacred history.
But when science has got its altar candle lit and
the ``mystic half-darkness'' in which the event is shrouded
begins to disperse, the congregation will surely
find ``the weight of the difficulties that press upon the whole
conception of the devil's personal reality, in
quite a peculiar way strengthened by the relation in which
he here steps forth toward Jesus. For as regards the devil,
the procedure in the whole difficult
struggle, the flatness in the arranging and introducing of the temptations
seems just as hard to reconcile with the
shrewdness and calculating cunning that were here supposed to
reach their climax, as, on Jesus' side, the patience
in continuing these negotiations, and in letting himself be led by the evil spirit
from place to place through the different stages
of the trial, is hard to reconcile with the ideal of
% --- p. 299 ---
\setcounter{footnote}{0}%
\opage{299}moral purity that must at once have withdrawn
from contact with the unclean''. Once it has really
dawned upon the enlightened congregation that the history of the temptation,
historically understood, is by its ``flatness'' equally compromising
for Christ and for the devil, then faith will probably
have to be content with the scientific explanation that
the whole ``is probably founded on a communication
from Jesus about something he had experienced in his soul in the
solitude of the desert, whereby it must for the rest be left
undetermined whether the disciples understood what was communicated
as something experienced inwardly or as something experienced outwardly''.

If one takes the reverend interpreter's
declaration, that it is scientific criticism that is to
test, find out and ground the historical truth of the
Gospel accounts, in full earnest, then one must, as to
the above-mentioned section, greatly wonder at our
clergy. For whether it be the Gospel of the temptation
or of the Annunciation or any other Gospel whatsoever
that constitutes the Text of the Day, the edifying
discourse proceeds without more ado from the presupposition
that the Gospel read aloud is in all parts historically
credible. The priest is silent about the difficulties, passes over
the critical misgivings and lets
nothing be noticed. It is, according to the interpreter's way of looking at things,
a manifest violation of the sacred truth.
Such a violation of the duty of truth cannot be excused
by ignorance; it must then have its ground
in this, that in the pulpit ``the unnatural'' is demanded
of the preacher: ``that he is here to keep his eyes tightly shut, or that, seeing, he is to renounce the ability to
see''.

% --- p. 300 ---
\setcounter{footnote}{0}%
\opage{300}It is true that the rule holds that one must not give
the simple-minded offense; but truth is after all always
truth, and the truth must out. If the principle of faith
is really one with the principle of knowledge, then there can
not possibly be one truth of faith in science and an
other in the congregation, one on the professorial chair and another in
the pulpit.

The nearest way to root out superstition and promote
religious cognition of truth would then be:
to exclude from the annual cycle of Gospels the sections
about whose historical credibility science entertains well-founded
doubt. But when the talk is of complete credibility
and historical reliability in all essential
features, for which sections of the sacred history does science dare
scientifically to vouch?

That ``the Transfiguration on the Mount'' and ``the Walking on
the Sea'' can come into consideration just as little as
``the history of the temptation'' or the ``fragmentary narratives''
that furnish ``the pre-history'', the interpreter himself
has already indicated.

\emph{The Transfiguration on the Mount} (Mtth. 17, Mark. 9,
Luk. 9) \emph{drops out;} it is not a picture with distinct
features of historical reality that we have before us here, but
``a mist-picture in mystic half-light''. He (Christ)
\emph{was transfigured before the disciples,} but in what the transformation
actually consisted criticism cannot find out.
``How the intimated metamorphosis is to be conceived:
as extending to the inner nature of the body --- and in
that case, how compatible with Jesus' true humanity
--- or as merely apparent with regard
to the outward appearance; --- how the presence
% --- p. 301 ---
\setcounter{footnote}{0}%
\opage{301}is to be conceived of the two heroes of Judaism''
(Moses and Elijah) ``the representatives of the Law and the Prophets:
real or visionary; --- what purpose is
to be ascribed to the described meeting, and what is
to be thought the subject of the conversation, --- how
the voice from the cloud and the words uttered are to be explained: all
this is a matter of questions to which a tolerably satisfactory
answer would have to be found, in order that the story told
might become an object of conception
and presentation. But from the nature of the narrative
it is evident that and why no such answer
can be given. For the three eyewitnesses'' (Peter, James
and John) ``were heavy with sleep, they did not know
what they said, for they were sore afraid''; when they
at last came to themselves and ``looked round about, they saw no man
save Jesus only''.

\emph{The Walking on the Sea} (Mtth. 14, Mark. 6)
\emph{drops out.} The interpreter maintains with philological right
that there is here no talk of a walking \emph{by} the sea or
\emph{toward} the sea, but literally of a walking \emph{on} the sea,
\emph{across the sea,} and refers to Job 9:8, where it is
said of God that he walks upon the deep as upon a
floor. The whole narrative, he says, ``obviously
has this explanation as its presupposition; for while the disciples
in the ship were in the midst of the sea, it is said that
Jesus would have passed them by, and they were amazed, supposing
that it was a spirit they saw before them''. So far
criticism, on the presupposition of the possibility of miracle,
makes no scientific objection to the event;
but --- naturam furca --- now it comes. The story
of Peter, who also wants to try to walk upon
% --- p. 302 ---
\setcounter{footnote}{0}%
\opage{302}the water, a characteristic scene, which completes the picture
and gives it a rich symbolic meaning, criticism cannot
approve. The edifying discourse, however, speaks
in an entirely different tone. ``These are written, that
ye might believe that Jesus is the Christ, the Son of God;
this word we would apply to our sacred narrative
read aloud, of the Lord's walking on the Lake of Gennesaret\dots{}
If we believe in him who has raised the dead, who has ruled
over the law of death, should we then not also believe in
him who has ruled over the law of gravity? We should
then also learn this from our Gospel, that what
matters is to have steadfastness, constancy
in faith, that we not, like Peter, begin in faith
the walk, but when we see the strong wind,
become afraid, doubt and sink''. That is very edifying indeed.
But when the identity of principle of faith and knowledge
is held fast, the foundation of this edifying talk becomes
hollow and floating. The preacher speaks of something that never
happened. Criticism has ``well-founded conjecture''
that ``the episode of Peter'' is a later addition, owing
to the Greek translation or revision of
the Gospel of Matthew --- ``a conjecture which is corroborated
by the whole character of this narrative, which reminds
one of the embellishing and intensifying tendency in the apocryphal
Gospels''.

``\emph{The Stilling of the Storm on the Sea of Galilee}''
(Mtth. 8, Mark. 4, Luk. 8) \emph{drops out, in order that the congregation}
may not take offense at the striking misunderstanding
that has crept into the presentation, a
misunderstanding so great that scientific
interpretation completely transforms the historical
% --- p. 303 ---
\setcounter{footnote}{0}%
\opage{303}character of the event. All three Gospels, namely, place the emphasis
unanimously on this, that storm and waves must
obey Christ's almighty word. But science understands
the matter differently: the whole significance lies in the contrast
between the Master's ``exalted calm'' and the disciples' ``anxious
timidity'', so that Jesus' threatening words ``were rather
probably addressed to the disciples than to the elements,
to weather and wind''. --- There are, if we follow the
scientific criticism, still many sections that on
account of historical unreliability must be excluded, indeed
so many that one cannot suppress the question
whether in the end a single one remains.

It is, we repeat it, by no means our opinion that
the reverend interpreter in his investigations has gone
too near the historical facts. On the contrary! if
the demand of science is to be satisfied in any measure,   % print: Fyldelst (as printed)
no one will surely be able to treat the evangelical
history with more caution or greater readiness
to let the wonderful step forth with the stamp of historical reality
everywhere where it is possible,
than has been done in the work here under discussion. But precisely
this endeavor: \emph{from science to secure
and ground the supernatural in the content of faith},
has, as we shall presently see, brought about a weakening
of the principles, a groping and wavering in the judgment
of the phenomena, which approaches the
unprincipled.

To meet the objections that the unbelieving criticism
(Strauss, Weisse, B. Bauer) has raised against the healing
of the centurion's servant (Mtth. 8, Luk. 7) as a
``healing at a distance, without there being any question of the
% --- p. 304 ---
\setcounter{footnote}{0}%
\opage{304}sick person's own faith'', the author remarks: ``although it
must be admitted that thought finds it easiest to appropriate
and hold fast the miracle where it appears bound
to faith as a middle term\dots{}, yet the
same evangelical miracle-type appears everywhere where the holy
will reveals itself as the mighty will, in the
bodily world as in the moral''. Here, then, the cause,
the holy, mighty will, is made the mark of the miracle.

But this mark is effaced again where the talk
is of the healing of the woman with the issue of blood (Mtth. 9,
Mark. 5, Luk. 8). As the narrative is given
in Mark and Luke, ``it leads namely with necessity
to the conception of Jesus' miraculous power
as working\dots{} without the ground of the activity being
to be sought in the determining might of the will\dots{} The healing
power is thus here presented as set in motion
by the mere outward touch, and the healing that followed
proceeds as a physical phenomenon,
where the thought involuntarily presses upon one of the electric
apparatus, filled with the strong power, which needs only
to be brought near a body in order to stream
out into visible effects''. Strictly speaking, such a
miracle ought now to be rejected outright; but one can instead
reject the mark, give up the definite purpose of the
mighty will, hold to the believer and
make faith the electric conductor of the miracle-power.

This scientific conception has edifying warrant
in Schleiermacher's sermons. ``Deeply and ingeniously'',
says the Interpreter, ``Schleiermacher has expressed himself on this
problem: that just as the spiritual influence that
% --- p. 305 ---
\setcounter{footnote}{0}%
\opage{305}proceeds from the individual has only a small part of its
ground in the determinate will and purpose, but such a
circulation of the spiritual forces takes place that the
best can be effected without the one who effects it being able to
say anything of it, so it was also with Jesus' miraculous
deeds in the bodily sphere, which just thereby removed
themselves from the usual course of things and approached
the characteristics of the spiritual life; there as well as here
faith, longing and yearning were the only condition
for becoming a partaker of the miraculous working''.

If the reader will not rest content with the Schleiermacherian
analogy, but finds the event unreasonable,
then critical interpretation still has the natural
way out of seeking the ground of the unreasonable in a misunderstanding
on the part of the evangelists. Matthew gives the narrative
briefly and places the power of the healing in Jesus' words:
\emph{Daughter, be of good comfort, thy faith hath made thee}
whole. One can now assume that this is precisely the historical element,
while the narrative in Mark and Luke ``makes
the impression of traditional elaborations, formed by reflection
on the manner and way of the healing''.

The natural way out is, however, the most dangerous. It is true that
at first glance it seems quite plausible that an
evangelist in all good faith may have come to
use one legend or another in which the miraculous, through additions
and ``traditional elaborations'', has been misunderstood:
it is then not the evangelist, but the legend, that
is to blame for the misunderstanding. But if we look more closely,
the matter stands differently. For the evangelist
who can take up legends that put the miracle in a false
light thereby precisely reveals that he lacks
% --- p. 306 ---
\setcounter{footnote}{0}%
\opage{306}view of the distinguishing mark of the miracle, that, to put it bluntly, he is not
able to distinguish between true and false miracles.

That an Evangelist may have erred in giving some historical outer side,
a name, a determination of time
and the like, is unessential, for an Evangelist is
no scientific, critical investigator of history; nor
is there anything exactly unreasonable in his having, by making use of reports
of legendary, uncertain origin, in his presentation
confused ``what was experienced inwardly with what was experienced
outwardly,'' for in the religious sense the outward
always receives its significance through the inward. But to suppose
of an Evangelist that in thoughtless credulity he has
approved narratives that in reality distort
the miracle, is by no means admissible, for that is to destroy
his credibility as an Evangelist. What
the interpreter adduces in this respect in excuse
of what is little credible in ``the embellishment of the miracle with
the woman with the issue of blood, with the walking on the sea, with
the stater in the fish's mouth, with the resurrection of ``the Jerusalem
dead, etc., is not an excuse but an accusation;
it is sheer doubt of the sacred authors' ability
to handle the historical material, veiled objections
to their credibility.

``We ought,'' says the reverend author, ``to let ourselves be guided by Scripture
itself to the true concept of the
holy word of revelation.'' Yes, so we ought. But this
principle is surely forgotten or pushed aside, every
time we demand a quite different presentation of the
wonderful from the one that is notoriously given to us; every
time we refer to ``later, askew, misinterpreting reflections''
and are on the point of setting our own so-called
% --- p. 307 ---
\setcounter{footnote}{0}%
\opage{307}scientific conception of miracle higher than the evangelical one.

How little, as far as the miracle is concerned, one can be minded
to let oneself be ``guided by Scripture to the true concept,''
the criticism of the miracle of the feeding
(Mark. 6, Mtth. 14, Luk. 9) offers a telling example. ``For
exegetical consideration,'' says the author, ``the miracle
that five thousand were fed with five loaves and two
fishes has as complete an outward warrant as any of
the evangelical miraculous deeds.'' None the less
scientific criticism must feel itself highly dissatisfied
with this miracle. The reason is ``that what
here confronts us is an act of power, in the absolute incomprehensibility
of the creating divine force\dots{}. A prodigy
(τέρας) stands before us that closes itself to every
effort of thought and dismisses every question that seeks
an answer as unanswerable.'' --- This uncertain conception of the supernatural, at odds
with itself, this wavering
reasoning for and against, these irresolute estimates
of reasonable, little reasonable and unreasonable events,
this unfruitful weighing of probable, little
probable and improbable miracles, is equally
ill-fated for religion and science: it makes
faith doubting and knowledge unprincipled.

It can be proved from the work of interpretation we are considering
that the sacred history must in principle exclude
science. By admitting the possibility of miracle
knowledge has set its absolute boundary and thereby spoken its last
word. History has, however, not a general possibility
but real facts, these signs, these
visions, these revelations, these miraculous deeds conceived in determinate
% --- p. 308 ---
\setcounter{footnote}{0}%
\opage{308}representations, as its content. If
science were in any way able to pass
a well-grounded judgment on the individual miracles
appearing in determinate phenomena, then it would first have to
discover the laws according to which the supernatural can
express itself and intervene in the natural. If it is now, as we
\emph{all know, wholly ignorant of these laws, then its}
criticism is without a standard: either it must let the miraculous deeds
stand for what they are: improbable, incomprehensible,
or it must seek to dissolve every single determinate
representation of every single determinate miracle into the indeterminate, indeterminable possibility
of ``the wondrous.'' In the
first case science will withdraw from the
sacred history, in the latter case its criticism will dissolve
the sacred history.

Our reverend interpreter has perhaps not reflected
that the real result of his work, so significant
for science, is a consummated
dissolution of the evangelical history. And yet the dissolution is
undeniable. There prevails in this interpretation
from first to last such a divorce between the
natural and the supernatural that only the natural
becomes historical, while the supernatural --- the in itself
indeterminate and indeterminable --- withdraws from every
possible historical apprehension. It is said in the introduction
to the first section that scientific criticism will ``secure
the historical kernel'' by separating it from ``the
embellishing setting'' and in particular ``maintain the historical
reality of a wondrous ground-fact in
Christ's entry into the world''; but whoever trusts
that here, with the assistance of criticism, he will find either
% --- p. 309 ---
\setcounter{footnote}{0}%
\opage{309}proofs for the reality of the wondrous ground-fact or
closer information about its historical character, will be
greatly disappointed.

\emph{How did it actually come about with the announcement
of the birth of John the Baptist?} One
does not know. Luke (1:5—25) describes it, to be sure,
in detail; by adducing the nearer circumstances and
incidental circumstances the Evangelist doubtless thinks that he makes clear
to us the conception of what is supernatural in the event and
tells a truthful story. But the interpreter is of
another opinion. He thinks that what the Evangelist here
tells is not any real event, not any
truthful story, but an Old Testament embellishment
of a family legend. ``Just as it is probable
that traditional legends were used that were preserved
in John's family, so it is obvious that the
Old Test.'s narratives of the birth of Samson and of Samuel
(Judg. 13. 1 Sam. 1) have served as models
in the individual features.'' What is it, then, that has actually
happened? One does not know: the event itself falls
outside history.

\emph{How did it come about with the announcement
of Christ's birth?} One does not know. Matthew
(1:18—25) and Luke (1:26—38) think, to be sure, that they are able to
impart reliable information and give us each
his contribution; but that is a pious fancy. Criticism
proves that there is no connection between ``the two
accounts,'' since the one obviously makes the other
superfluous. ``Nothing is surely in itself more natural than
that, about an event in which the supernatural lost
itself in mysterious obscurity, there should be different
% --- p. 310 ---
\setcounter{footnote}{0}%
\opage{310}narratives in circulation; and in any case there has, on the part of the narrators,
been no thought of uniting them
with one another; for it cannot be denied that --- after the
angelic revelation to Mary --- the dream-vision for Joseph
stands little motivated.''

\emph{Can the angels' revelation to the shepherds}
(Luk. 2:8—20) \emph{be regarded as a historical event}?
Yes, in the pulpit, but not in science!
Here, to be sure, is a ground-fact; but when criticism has
got ``the setting'' separated from ``the kernel,'' then
the wondrous ground-fact disappears in the obscurity of legend
and finally dissipates in an ingenious psychological ambiguity.
For we must, with Schleiermacher, ``assume that this
history, only after Jesus's reputation had been established,
was drawn forth from oblivion by the recollection of individuals;
but that it is neither an air-bubble puffed together out of nothing
nor a deliberate invention is vouched
for us by the ease which the Christians of that region had in
questioning Mary, who necessarily must have known about
it, or those disciples with whom she had most
lived, about the nature of the matter.'' Thus something must
have happened, something wondrous, but what sort of thing is it?
One does not know. The purpose of the event is
namely not to be sought in ``the more remote effect which it
produced, after its resurrection, so to speak, but
only in the immediate one on the shepherds themselves. The private revelation
to the shepherds can now be explained in a quite
natural and yet almost supernatural way by this, that someone
or other among them has, ``in his peculiar piety,''
asked for a sign. ``This was granted them, as history
teaches, and in this, in such great times not rare,
% --- p. 311 ---
\setcounter{footnote}{0}%
\opage{311}satisfaction also of the longing of individuals lies the remarkable and divine.''

\emph{In his public life Christ himself is the wondrous
ground-fact;} but that criticism cannot here either
come to terms with the historical phenomena
of this ground-fact we have already seen in
many examples. In order not to reject all individual
miraculous deeds, and yet to keep one's head above water with regard to the multiplicity,
the reasoning understanding ascribes to and denies the Savior natural and supernatural gifts, as
it may best happen to fall out. That Jesus, according to what we have
heard about the centurion's servant, did have gifts
to heal a sick man at a distance, is tacitly conceded;
but that he also had gifts to be able to
see an ass at a distance, is not conceded. ``The ass''
--- it is said of the entry into Jerusalem (Mtth. 21,
Mark. 11, Luk. 9) --- ``is fetched by the two disciples
from the village lying directly opposite, probably Bethphage.''
All three Evangelists are, to be sure, agreed
in laying stress on the wondrous prediction that
they would ``find'' the ass, and thus hint at a ``mysterious
connection in the whole occurrence.'' But in this
criticism cannot support the sacred authors.
Science must find it probable that ``either
the necessary further information was added by
Jesus, or else a previously made arrangement on Jesus's part must
be understood.''

Yet another trait. The interpreter lets it be heard
that the Transfiguration on the mount is hard to reconcile with
``Jesus's true humanity''; here, then, no supernatural
ground-fact is to be thought of; but can
% --- p. 312 ---
\setcounter{footnote}{0}%
\opage{312}\emph{the Resurrection and the Ascension} then be better reconciled
with ``Jesus's true humanity''? Here we must hold
fast to a mysterious fact. But the mysterious
fact is so far from clearing itself up in the presentation
that it on the contrary makes the presentation unclear. ``The mutual
divergences with which the four evangelical accounts
of the resurrection cross one another'' have
so disturbing an effect that ``only with great uncertainty
can a more or less coherent narrative be produced,
by putting them together with one another'';
in other words: what was really to be told, the supernatural
\emph{How} of the occurrences, is not told;
what was here to be attested is not attested.
And that is quite in order, it cannot be otherwise.
The ground of the uncertainty lies by no means in
any carelessness or neglect on the part of the eyewitnesses,
it lies in the nature of the event, in ``the overpowering
impression by which the first eyewitnesses had to be seized,
in which connection it comes especially into consideration that the tradition
in its individual features has here had its first warrant in
the uncertain observation of the female witnesses.''

As it goes with the Gospel of the Resurrection, so it goes
also with \emph{that of the Ascension.} The Ascension is attested
as a ground-fact, as an article in a creed,
but of the Ascension itself no history is given.
And why not? Because the wonderful
ground-fact stands in manifest conflict with all history.
What cannot appear in a determinate form in
a determinate manner, cannot either occur at a determinate
place at a determinate time, cannot become event;
but what by its nature cannot become event,
% --- p. 313 ---
\setcounter{footnote}{0}%
\opage{313}cannot become historical either. Criticism sets an
insurmountable gulf between the fact of the Ascension and
``the \emph{form and manner} in which this fact, this
wondrous entrance into heavenly glory, is to be
thought accomplished.'' Yes, the Ascension is, when we look
more closely, not even a proper fact; it is
a ``closed mystery,'' but a closed mystery is
no fact. ``Had there not,'' with this the reverend
interpreter concludes, ``been here a closed
mystery already for the Apostles, then neither the omission
of the greatest σημεῖον'' (sign) ``in the two
apostolic Gospels would be explicable, nor the obscurity
in which the allusion in the two other Gospels
is held.'' This is the end of the sacred history
--- in science.

This remarkable circumstance, that so clear, so acute an author,
so practiced in critical thinking, so familiar with his great scientific
apparatus, as Dr. H.
N. Clausen, has by carrying through his principles
arrived at a result that is precisely the opposite of
that at which he originally aimed, raises the presumption
that in the principles themselves there must lie hidden
a bewitching obscurity, a dazzling ambiguity.
We shall endeavor in brief to uncover this
ambiguity.

That, according to principle (1), a distinction must be made
between profane and sacred history is undeniable;
but the question is wherein the distinction consists and how
deep it is grounded in the matter. The ``historical credibility''
is of another nature, of another origin, and has a   % print: Natnr (as printed)
quite different ground in sacred history than in profane.

% --- p. 314 ---
\setcounter{footnote}{0}%
\opage{314}To the concept of \emph{credibility} corresponds the concept of faith;
the credible is the object of faith. From this it follows that
the concept of credibility must change when the concept of faith
changes. But now there is, at the evangelical-Protestant
stage of cultivation, hardly anyone, I should think, who will
deny that religious faith (fides religiosa) is essentially
different from what is called in the common sense
historical faith (fides historica).

Merely historical faith concerns events that
in their time were the object of outward, immediate
experience, and about which the eyewitness can to that extent
be assumed to have had reliable knowledge. Events
of this nature all belong under the law-governed
order of things that forms the firm foundation of outgoing observations,
of sensory experiences. In the
apprehension of such events the outer will everywhere
be determinative of the inner. The eyewitness has to
keep to what has shown itself, to reproduce what was experienced,
as a mirror gives the image of the object, not
to mix in anything of his own, not to put color on
the facts. By meeting these demands
an eyewitness is credible. He, on the other hand, who hears
or reads the account and thus apprehends the event
at second hand, cannot, like the eyewitness,
have an immediate knowledge of what was experienced, but must
content himself with a faith, a historical faith.

If historical faith now rested solely
on an unbounded trust in the individual veracity of the narrating
eyewitness, such a faith could not possibly
support any scientific edifice. History
would on these terms not be a science.
% --- p. 315 ---
\setcounter{footnote}{0}%
\opage{315}But precisely the circumstance that here there is talk only of
facts that are subject to natural conditions and to the
universally valid laws of spirit of human life, makes it
possible that historical faith can be grounded, in that the accounts
are subjected to a criticism, and their credibility
is tested.

The credibility of a historical event is sufficiently
grounded when it has been made evident that
one cannot deny such an event without
denying all history. Probability is the measure of credibility:
the more probable an occurrence is, the more
credible. But probability is in turn determined by
the inner as well as outer connection of the facts adduced. Thus an occurrence, an action, a
trait of character, which, torn loose from connection with the
rest, would be improbable, can by all-sided, comprehensive
comparison become probable. Grounds of probability
of tested validity are sufficient for profane
history.

Otherwise in sacred history. Here the grounds of probability must
be reduced to a subordinate value;
for that on which everything depends is the miracle, and the miracle
is by its nature the improbable. The difference
between a miraculous deed and a plain historical fact
is so thoroughgoing that historical faith
transforms itself and merges into the religious. It is here,
as far as apprehension is concerned, not the outer that immediately
determines the inner, it is conversely the inner
that essentially determines the outer. Consider whichever
of the evangelical miracles one will, then it shows itself
at once that the phenomenon, the outer, admits of a double
% --- p. 316 ---
\setcounter{footnote}{0}%
\opage{316}explanation.
It may be a disappointment, a deception, a strange
coincidence, a surprising natural event,
and it may be a work of God. Christ raises
the dead: the daughter of Jairus (Mtth. 9), the widow's son at Nain
(Luk. 7), Lazarus of Bethany (John 11). Of the first two
raisings it is now at once evident that they could be explained by the
eyewitnesses in two ways: the unbelievers
might indeed say that the so-called dead were only
apparently dead. As for the raising of Lazarus, the
palpability with which the miracle appears has not
been able to prevent a part of those present
from \emph{having regarded the whole as a deception.} It is expressly said
that some of them went to the Pharisees and
informed against Jesus, and that the Pharisees and the chief priests, precisely
on account of this information, resolved in council that
``the seducer of the people'' should die. If we consider the miracle
of the resurrection\footnote[1]{Cf. \emph{On Obstacles to and Conditions for Spiritual Life in the Present Age.} Postscript, pp.\ 297—310.}, the same thing shows itself. ``There was
(Mtth. 28) a great earthquake; for the angel of the Lord descended
from heaven, came and rolled back the stone from
the entrance and sat upon it. But his appearance was
like lightning, and his raiment white as snow. And
the guards trembled for fear of him and became as
dead.'' --- Immediately after these words one would now suppose
that the soldiers on watch really had seen
the angel in the lightning and thus had become eyewitnesses of
a real miracle. But that this was not the case
is clearly evident from this, that the guard let itself be bribed
by the chief priests. Had those on watch
% --- p. 317 ---
\setcounter{footnote}{0}%
\opage{317}regarded the whole as anything other than a terrible natural event,
they would surely not have had the courage to lie.

If it now lies in the nature of the case that the apprehension of
miraculous events is conditioned by the eyewitness's religious
faith\footnote[1]{On the logical impossibility of apprehending a miracle as any other plain historical fact, cf.\ The Logic of the Fundamental Ideas II, pp.\ 208—210.}, then the Evangelists too cannot
confine themselves to reporting what they have plainly
heard and seen, but must present and describe what they
\emph{believed they had heard and seen,} what they really experienced
in faith. The sacred history has no foundation of knowledge;
it is a history of what was experienced in faith,
a history of faith, a history of believers for believers.

From this point of view the ambiguity contained in the interpreter's
principles will already begin to
show itself. The reverend author emphasizes that
the Evangelists' relation to the material that ``fills out the
evangelical account'' cannot have ``the same definite
firmness and determinateness as where the historian is
placed before a historical material in which nothing goes beyond
the boundary of the customary and in itself credible'';
he maintains, in good agreement with principles
(2) and (3), that where the historical
content is of so peculiar a character that ``the
whole moves on the dark, holy ground of mystery,''
there ``the religious way of looking cannot
remain without determining influence on the critical
assessment of the historical character of the Gospels.'' Here
is the ambiguity. The ambiguous lies in the way
in which the religious is combined with the scientific.

% --- p. 318 ---
\setcounter{footnote}{0}%
\opage{318}By calling religious faith to his aid, the interpreter
thinks he has won a sure foothold against those
who, in apprehending the miracle as ``the plainly contrary to nature,
that which is incompatible with the laws of thinking as with those of nature,''
would maintain ``that the evangelical
history cannot have taken place in the way it
must have, according to the common understanding of the evangelical
account.'' Here it is now
at once striking that the author is in the main not really
at odds with the opponents he wants with all his might to combat.

The critics who resolutely deny the possibility of miracle
(Strauss, Weisse and others) arrive with full scientific
consistency at the well-known result:
``the divine cannot have happened in this way; what
happened in this way cannot be divine.'' But our
interpreter, who has declared himself at odds with the modern thinkers
as to the presupposition, has, strangely enough, become
agreed, or at least almost partly agreed, with them in the consequence.
By holding fast to the separation of what he calls
``the outer side of the Gospels and the evangelical core-essence''
his interpretation is, as we have seen above,
in all essential points a veiled confirmation
of the Straussian proposition: ``the divine cannot
have happened in this way.'' The author promises himself something
great from the separation he invokes; by it ``it will
become possible for apologetics victoriously to repel the unwarranted
conclusions which anti-evangelical criticism will
draw from it.'' But it will not become possible; anti-evangelical
criticism will win precisely on account of the separation;
it will ask, and with full right: ``what good
% --- p. 319 ---
\setcounter{footnote}{0}%
\opage{319}does it do that something wonderful is assumed to have happened here,
when one cannot be told what it is
that has happened?''

If it is to be more than a figure of speech that the sacred
history is a history of revelation, then it is
really not admissible to let the supernatural hide in
a ``mystic half-darkness'' and keep itself as ``a closed
mystery'' behind the legendary accounts. To say
of the Evangelists' relation to the given material, the material in
their narratives, that this relation cannot have ``the same
definite firmness and determinateness'' as that of other historians,
is, bluntly, to admit that the evangelical
history, strictly speaking, lacks credibility, that the history of revelation
is a history that reveals nothing, ---
lucus a non lucendo.

As a history of revelation the sacred history
differs from the profane in this, that it has \emph{religious}
faith as its \emph{original and independent principle of cognition.} Revelation and faith are correlative concepts.
As revelation is, so is faith, and as
faith is, so is revelation; their essentially operative
unity is the spirit; but the revealing expression of the spirit
is the word. Let us then once more fix our gaze on the supernatural
facts, to see what lies in them.

The miracle by its nature cannot possibly be any
plain outward object of experience. The sensory
eye sees the lightning from heaven, but the eye of faith sees
the angels at the grave. The first immediate impression
is terror; out of terror understanding is born,
and out of understanding the presentation in the word. It is
the work of the spirit, and as the spirit is, so is the word.
% --- p. 320 ---
\setcounter{footnote}{0}%
\opage{320}If the word is legend-like and floating, without ``firmness and determinateness,''
then the spirit is without certainty, without clarity,
without truth; but if there is really truth,
historical truth in the working of the spirit, then there is also
really truth, historical truth in the word.

The ambiguity of the principles now shows itself from a
new side. The reverend interpreter feels full well that
the merely scientific principle, which is otherwise sufficient
when it is a matter of a criticism of a piece of world history,
is by no means sufficient when it is a matter of
a criticism of the sacred history. So he calls
faith to his aid; but in all indeterminateness, for ``the religious
way of looking,'' which ``cannot remain without
determining influence on the critical assessment of
the historical character of the Gospels,'' is floating and indeterminate.
The scientific is the main thing, faith only
an emergency aid, which is to be trimmed to science.

By this, albeit cautious, yet most precarious
combining of faith and knowledge the interpreter has prepared
for himself a great embarrassment. On the one side
he combats the doctrine of myth, ``the mythicizing treatment,''
most vigorously. ``An exertion to
prove that what the Gospels contain and hand down
is not an unconscious or half-conscious pious work of imagination,
but an account of a real life-history,
real occurrences and deeds,'' must, he says,
``already now be regarded as untimely.'' Why? Because the way of looking
is so unreasonable that the unreasonable is evident of
itself. That is indeed gratifying for faith. But the gratification
is not particularly great; for on the other side he combats
just as eagerly the old-church way of looking,
% --- p. 321 ---
\setcounter{footnote}{0}%
\opage{321}that we have before us everywhere in the Evangelists' accounts ``a
real life-history, real occurrences and deeds.'' It is a necessary concession that
must be made to science.

If the reverend author could now only in
some way see his way to reconcile the two mutually
conflicting ways of looking with each other: then such
a reconciliation would be an eloquent proof of the unity of faith
with knowledge. But the reconciliation does not come about. ``At the
same time,'' he says, ``as such strong and
such well-founded complaints are raised against the arbitrariness and injustice
to which the sacred history is delivered up
by the most recent criticism --- with all apparent strictness
in the scientific form, yet in a little scientific
spirit --- the serious questions cannot be pushed
aside: whether criticism has not been tempted and egged on to
so hostile a bearing by an old, often repeated and
very pressing challenge from the side of apologetics?
whether the apologetic task can be rightly posed\dots{},
when it'' (Apologetics) ``continues, as hitherto, to let   % print: Apologotiken (as printed)
faith in the Bible step forth with peremptory language and anathema against
every concession of phenomena that point to a
human side of Scripture, continues as hitherto to
demand an unconditional denial of what every comparative
and testing consideration forces upon us as
indubitable certainty?'' In emphasizing these difficulties
the reverend author is perfectly right, and the criticism of the Gospels
will yet, after it has long since abandoned his
standpoint, with renewed interest return to his
instructive work and always count it to him as true
% --- p. 322 ---
\setcounter{footnote}{0}%
\opage{322}merit that he has honestly counteracted exaggerations
and kept a watchful eye on the difficulties.

But that no interpretation, were it ever the most learned,
talented and acute, is able from the standpoint
here chosen to remove the difficulties pointed out
is clear as day. What good does it do
to appeal to a ``separation of the essential
and the unessential,'' when the separation, according to what we have
learned from the interpreter's own words, expressly shows that the
essential is ``a closed mystery,'' which cannot
assume historical form, and the account in its historical
form ``the mythicizing garb,'' ``the embellishing
setting,'' which goes around ``the kernel,'' the
unessential.

Only on the condition that the principle of \emph{faith, and thereby}
of revelation is allowed to count for
what it is in truth, \emph{the sole decisive
principle of cognition of the sacred history,} can
the tangled knots be untied. For herein lies: 1) that
the elements in the evangelical history are both \emph{outward}
--- stemming from all that in outer
reality was seen and heard and experienced by eyewitnesses ---
\emph{and inward,} stemming from the individuals' particular apprehension
of what they have in faith heard and seen and deeply
lived through; 2) that the outward and inward elements are
so fused together, worked together and assimilated
in the sacred witnesses; \emph{that no criticism is}
in a \emph{position to separate them from one another
in the final sense;} 3) that the historical, in the events
themselves as in the corresponding accounts, is neither
a merely outward nor a merely inward
% --- p. 323 ---
\setcounter{footnote}{0}%
\opage{323}reality, but the unity of spirit and truth, of the
outward determined by its inward, thus \emph{Spiritual Reality,
reality for the spirit.}

Under these presuppositions the difficulties can be
removed, for the mistaken view from which they spring can be precisely
pointed out. The Holy Scripture has, as our reverend   % print: paayises (as printed)
interpreter particularly insists, not only its
divine but also its human side\footnote[1]{\emph{On Fundamental Determinations in the Self of the Spirit,} cf.\ Philosophy of Religion, pp.\ 458—486.}. But
apologetics and negative criticism are in this respect,
in the midst of all their disagreement, precisely agreed in proceeding
from the same mistaken view.

Negative criticism misjudges the spirit, goes behind
the word, tests the outward facts and discovers
such a multiplicity of contradictions and inner
discrepancies that the whole must seem incredible.
It sets a selfless idea in the place of the spirit and makes the word
a myth: here the human has smothered the divine.

Orthodox apologetics nails its dogma fast
and makes the evangelical history an infallible work
of God's infallible Holy Spirit. But in favor of its
assertion it misjudges the spirit, stares at ``the outer,''
goes behind the word and hangs on to the letter. All
the events are to be, down to the smallest features, outward,
rendered with so outward, photographic
exactness that the Holy Spirit must have dictated everything
down to jot and letter: here the divine has
smothered the human.

% --- p. 324 ---
\setcounter{footnote}{0}%
\opage{324}Then comes the apologetics of mediation. It will reconcile
the human with the divine, and hastens
to commit the same error as its opponents.
Instead of seeking the union of the divine and
the human in the spirit through the word, it dissects
the word, keeps to sentences, overlooks the spirit,
the inward, and in the most naive way sets
about the outward. It seeks the actual
fact outside the word; it peeps and spies,
whether it might not here or there catch a glimpse of what the
real, the wondrous, actually is that stands behind
the account. But all its spying is in vain.
The real itself, the wondrous, will not let itself
be seen; it is and remains an indeterminable ``ground-fact,''
``a closed mystery'': the apologetics of mediation
discovers nothing.

For fear of ``dualism'' the mediating apologetics
has entangled itself in the most dangerous dualism of all, the dualism
of ambiguities: \emph{one truth in the scientific,
another in the edifying; one Gospel
at the lectern, another in the pulpit};
\emph{one faith within knowledge, another outside}
knowledge. Should it not be time that we got this
dualism cleared up. The ambiguities must go! It
is not merely a pious wish, a quiet sigh of the heart in the name
of faith; it is the demand of enlightenment, the stern command of the
wrathful reason, of the spirit of the age in its dark majesty.

``\emph{Cardinal, ich habe das Meinige gethan.
Thun Sie das Ihre}.''
\signature{R. \emph{Nielsen}.}

\end{document}
